% !TEX root = ../main.tex
\section{Conclusiones}

El sistema difuso Mamdani propuesto modela el riesgo académico como una transición gradual entre conjuntos lingüísticos. La combinación de funciones de pertenencia triangulares y trapezoidales, diez reglas IF-THEN, agregación máximo y defuzzificación centroide produce un resultado numérico explicable y reproducible, sin recurrir a aprendizaje automático ni estadística predictiva.

El análisis de sensibilidad demuestra que el sistema discrimina entre escenarios contrastantes con valores claramente diferenciados y sin saltos artificiales. La sincronización entre motor TypeScript, reporte LaTeX y presentación PptxGenJS asegura que los tres entregables compartan una única definición del sistema, eliminando el riesgo de divergencia entre documentación y código.

\subsection{Aportes del proyecto}

\begin{itemize}
  \item Motor determinista en TypeScript puro auditable línea por línea.
  \item Reporte académico con figuras vectoriales generadas a partir del motor.
  \item Presentación editable PptxGenJS que se regenera con el caso analizado.
  \item Interfaz web interactiva con descargas integradas de los tres entregables.
\end{itemize}

\subsection{Limitaciones}

\begin{itemize}
  \item Las funciones de pertenencia se calibran por criterio experto. Una revisión periódica con docentes asegura su validez en distintos contextos académicos.
  \item La base de reglas crece de forma combinatoria con el número de variables; sistemas más complejos requerirán jerarquías o metareglas.
  \item El sistema produce alertas, no diagnósticos: la decisión final requiere validación humana por el tutor académico.
\end{itemize}

\subsection{Trabajo futuro}

Como continuación natural se contemplan: tuning automático de los conjuntos mediante ANFIS preservando la interpretabilidad, integración con plataformas LMS para alimentar las variables, evaluación longitudinal por semestre y comparación con expertos humanos. La extensión a sistemas tipo-2 \citep{castillo2008type2} permitirá modelar la incertidumbre en los propios parámetros de las funciones de pertenencia.

\section{Checklist de validación}

\begin{itemize}
  \item Variables de entrada y salida definidas con rangos y conjuntos.
  \item Funciones triangulares y trapezoidales graficadas a partir del motor.
  \item Reglas IF-THEN visibles con justificación lingüística.
  \item Inferencia Mamdani con AND mínimo verificada paso a paso.
  \item Agregación máximo y defuzzificación centroide documentadas.
  \item Resultado crisp y etiqueta lingüística reportados para cuatro casos.
  \item Análisis de sensibilidad con casos contrastantes incluido.
  \item Apéndice con código fuente del motor y activación completa de reglas.
  \item Interfaz interactiva con descargas integradas.
  \item Sin machine learning, redes neuronales, regresión ni modelos estadísticos.
\end{itemize}
